%========================================================================================
%  SINTESI DELLA TESI (versione corretta) — Lorenzo Caivano
%  Compila questo file da solo (è indipendente dalla tesi).
%========================================================================================
\documentclass[12pt]{report}

\usepackage[utf8]{inputenc}
\usepackage[T1]{fontenc}
\usepackage[a4paper, margin=1in]{geometry}
\usepackage{setspace}
\usepackage{parskip}
\usepackage{titlesec}
\usepackage{amsmath}
\usepackage{amssymb}
\usepackage{xcolor}
% Parti aggiunte da Irene: in grassetto e in blu
\newcommand{\irene}[1]{{\color{blue}\bfseries #1}}

% Titolo di capitolo semplice e centrato (niente "Capitolo 1")
\titleformat{\chapter}[display]
  {\normalfont\bfseries\centering}{}{0pt}{\Large}

%----------------------------------------------------------------------------------------
%  DATI DELLA TESI
%----------------------------------------------------------------------------------------
% Italiano
\newcommand{\cognomeIT}{Caivano}
\newcommand{\nomeIT}{Lorenzo}
\newcommand{\matricolaIT}{889490}
\newcommand{\titoloTesiIT}{Studio delle degenerazioni nell'analisi del rumore delle pulsar}
\newcommand{\relatoreIT}{Prof. Alberto Sesana}
\newcommand{\correlatoreIT}{Dott.ssa Irene Ferranti}
\newcommand{\dataSedutaIT}{23--25 Novembre 2026}
\newcommand{\corsoDiLaureaIT}{Fisica}
\newcommand{\recapitoTelefonicoIT}{+39 3391306307}

% Inglese
\newcommand{\surnameEN}{Caivano}
\newcommand{\nameEN}{Lorenzo}
\newcommand{\matricolaEN}{889490}
\newcommand{\thesisTitleEN}{Study of degeneracies in pulsar noise analysis}
\newcommand{\advisorEN}{Prof. Alberto Sesana}
\newcommand{\coadvisorEN}{Dr. Irene Ferranti}
\newcommand{\defenseDateEN}{23--25 November 2026}
\newcommand{\degreeCourseEN}{Physics}
\newcommand{\phoneEN}{+39 3391306307}

\begin{document}
\thispagestyle{empty}

%========================================================================================
%  SINTESI IN ITALIANO
%========================================================================================
\chapter*{Sintesi}
\begin{center}
\textbf{Nome e Cognome:} \nomeIT \, \cognomeIT \\[0.3cm]
\textbf{Matricola:} \matricolaIT \\[0.3cm]
\textbf{Titolo (trad.):} \titoloTesiIT \\[0.3cm]
\textbf{Relatore:} \relatoreIT \\[0.3cm]
\textbf{Correlatrice:} \correlatoreIT \\[0.3cm]
\textbf{Data della Prova Finale:} \dataSedutaIT \\[0.3cm]
\textbf{Corso di Laurea:} \corsoDiLaureaIT \\[0.3cm]
\textbf{Recapito Telefonico:} \recapitoTelefonicoIT \\
\end{center}

\vspace{1cm}
\sloppy

\noindent
Le \textit{millisecond pulsar} (MSP) sono stelle di neutroni in rapida rotazione che
emettono impulsi radio con una regolarità paragonabile a quella degli orologi
atomici. Un'onda gravitazionale che attraversa lo spazio tra la Terra e una pulsar
perturba il tempo di percorrenza degli impulsi, che arrivano leggermente in
anticipo o in ritardo rispetto a quanto previsto. Monitorando per decenni i
\textit{times of arrival} (TOA) degli impulsi di un insieme di pulsar, un
\textit{Pulsar Timing Array} (PTA), è quindi possibile osservare onde gravitazionali
con frequenze dell'ordine dei nanohertz ($10^{-9}$--$10^{-7}$\,Hz). \par

In questa banda la sorgente principale attesa è il fondo stocastico di onde
gravitazionali (\textit{Gravitational Wave Background}, GWB), prodotto dalla
sovrapposizione incoerente dei segnali emessi dalla popolazione di binarie di buchi
neri supermassicci. Il GWB si manifesta come un segnale a bassa frequenza comune a
tutte le pulsar e correlato tra di esse in modo caratteristico; nel 2023 diverse
collaborazioni, tra cui lo \textit{European Pulsar Timing Array} (EPTA), ne hanno
riportato una prima evidenza. \par

La rivelazione del GWB è tuttavia resa difficile dalle numerose sorgenti di rumore
presenti nei TOA degli impulsi: irregolarità nella rotazione della pulsar
(\textit{red noise}, RN), variazioni della densità di elettroni lungo la linea di
vista (\textit{dispersion measure}, DM) e diffusione degli impulsi nel mezzo
interstellare (\textit{scattering}). Questi processi producono nei residui di
timing un rumore ``rosso'', con più potenza alle basse frequenze, descritto dalla
stessa forma spettrale a legge di potenza del GWB. In una singola pulsar il
segnale gravitazionale è quindi difficilmente distinguibile dal rumore, e le
diverse componenti del modello di rumore possono risultare degeneri tra loro.
Un modello di rumore accurato è perciò essenziale, perché un GWB potrebbe essere
assorbito dal rumore della pulsar o, al contrario, imitato da esso. \par

Il mio lavoro di tesi si è concentrato sullo studio delle degenerazioni tra i
parametri del modello di rumore di una singola pulsar, PSR~J1640+2224, inclusa
nella seconda \textit{data release} dello \textit{European Pulsar Timing Array}
(EPTA DR2). Per lo studio sono stati analizzati sia i dati reali sia dati simulati
in modo realistico. \par

L'analisi è stata svolta su due versioni dei dati di J1640+2224: DR2new+, che
contiene solo le osservazioni EPTA acquisite con la strumentazione più recente,
e DR2low, che le combina con osservazioni a bassa frequenza (30--200\,MHz) dei
radiotelescopi LOFAR e NenuFAR. Il rumore bianco è stato fissato al
\textit{maximum likelihood value} ottenuto da analisi precedenti a questa tesi,
mentre i processi rossi sono descritti con sei parametri liberi, un'ampiezza e un
indice spettrale per ciascuno tra red noise, variazioni di DM e scattering.
L'inferenza bayesiana è stata condotta con Discovery, un nuovo codice basato su
JAX: il confronto tra i due dataset mostra che solo con DR2low, grazie alla
copertura in frequenza, i parametri del rumore cromatico risultano vincolati. Su
DR2low le posterior di Discovery coincidono con quelle ottenute con Enterprise, il
codice di riferimento della comunità, e mostrano una marcata degenerazione tra DM
e scattering. \irene{Questa degenerazione crea una \textit{posterior distribution}
bimodale nei parametri della DM, con un picco che contiene il \textit{maximum
likelihood value} ad ampiezza bassa ($\log_{10}A \sim -15$) e \textit{spectral
index} $\gamma \sim 4$, e un picco ad ampiezza più elevata e \textit{spectral
index} più basso. Inoltre genera una \textit{joint posterior} fra DM e scattering
con forma a ``L'', sintomo, appunto, del fatto che le due componenti non vengono
correttamente separate. Attraverso l'analisi di simulazioni realistiche,
l'obiettivo di questa tesi è capire se componenti di rumore non modellate, ma
presenti nei dati, possano generare questa degenerazione.} \par

\irene{Assumendo che un GWB sia presente nei dati, questo rappresenta un red noise
acromatico aggiuntivo rispetto a quello generato dalle irregolarità nella
rotazione della pulsar.} Per studiarne l'effetto ho simulato dati con le stesse
caratteristiche di J1640+2224, iniettando un fondo con $\gamma = 13/3$ e ampiezze
comprese tra $\log_{10}A = -15.5$ e $-14$ (compatibili con le attuali misure dei
PTA), e li ho analizzati con il modello a sei parametri (RN, DM e scattering)
privo del termine di GWB. Per $\log_{10}A = -14$ il segnale \irene{provoca una
degenerazione fra DM e scattering con caratteristiche simili a quelle osservate
nei dati, seppur con significatività molto inferiore.} \par

Ho inoltre incluso un termine di \textit{band noise}, un processo rosso che agisce
solo sui TOA della banda LOFAR. \irene{Analizzando simulazioni in cui il band noise è
iniettato, ma non modellato, ottengo che questo rumore genera un bias in tutti i
processi, ma soprattutto nella \textit{posterior distribution} della DM, che si
sposta verso ampiezze più alte e \textit{spectral index} più bassi rispetto ai
valori iniettati. La presenza di un band noise non modellato crea dunque un picco
ad ampiezze e \textit{spectral index} compatibili con il picco secondario delle
\textit{posterior distribution} ottenute con i dati reali.} Inoltre, le simulazioni
mostrano che, a causa dell'elevato rumore bianco dei dati LOFAR, questo processo è
rilevabile solo per ampiezze $\log_{10}A \gtrsim -12.5$; sui dati reali se ne
ottiene solo un limite superiore, \irene{simile a quello che si ottiene sui dati simulati
iniettando un band noise con $\log_{10}A \sim -12.5$}. Di conseguenza, la presenza
di band noise nei dati reali non può essere esclusa e la sua ricerca sarà oggetto
di studi futuri. \par

In conclusione, l'analisi di J1640+2224 mostra che Discovery riproduce i
risultati di Enterprise e che la stima del rumore di questa pulsar è dominata
dalla degenerazione tra le due componenti cromatiche, DM e scattering, che solo
una buona copertura in frequenza radio permette di ridurre. Le simulazioni
indicano inoltre che un GWB di ampiezza paragonabile a quella oggi osservata, se
non modellato, può causare una degenerazione fra DM e scattering, e che un band
noise non modellato può generare un forte bias nei parametri della DM. Questi
risultati sottolineano l'importanza di una modellizzazione dettagliata del rumore
delle singole pulsar, presupposto essenziale per la rivelazione del GWB con i PTA.
\par

%========================================================================================
%  ABSTRACT IN INGLESE
%========================================================================================
\newpage
\chapter*{Abstract}
\begin{center}
\textbf{Name and Surname:} \nameEN \, \surnameEN \\[0.3cm]
\textbf{Student ID:} \matricolaEN \\[0.3cm]
\textbf{Title:} \thesisTitleEN \\[0.3cm]
\textbf{Advisor:} \advisorEN \\[0.3cm]
\textbf{Co-advisor:} \coadvisorEN \\[0.3cm]
\textbf{Date of Dissertation:} \defenseDateEN \\[0.3cm]
\textbf{Degree Course:} \degreeCourseEN \\[0.3cm]
\textbf{Phone Number:} \phoneEN \\
\vspace{1cm}
\end{center}
\sloppy

\noindent
English abstract text. \par

\end{document}
